% =====================================================================
\section{\'Etica y manejo de datos}\label{sec:etica}
% =====================================================================
% Fundamento: LOPDP \cite{lopdp2021}, su Reglamento \cite{rlopdp2023} y
% el Codigo de Etica de la ACM \cite{acm2018}. Cite articulos solo si
% los ha leido en el texto oficial y dicen lo que usted afirma.

El tratamiento de los datos de los participantes se rige por la Ley Org\'anica de Protecci\'on de
Datos Personales \cite{lopdp2021} y su Reglamento \cite{rlopdp2023}; la conducta del equipo, por
el C\'odigo de \'Etica de la ACM \cite{acm2018}.

\subsection{Consentimiento informado}
Antes de cada sesi\'on, cada participante recibe y firma el modelo de consentimiento del
Anexo~\ref{anx:consentimiento}. Los consentimientos firmados no se publican en el repositorio.

\pendiente{Describir c\'omo se entrega, explica y archiva el consentimiento; qu\'e ocurre si un participante lo revoca.}

\subsection{Anonimizaci\'on}
% Ejemplo de procedimiento: codigos P01, P02... por participante; la
% tabla codigo-identidad vive solo en la carpeta privada; transcripciones
% sin nombres, cedulas, telefonos ni correos; voces/rostros no se suben.
\pendiente{Procedimiento paso a paso de anonimizaci\'on aplicado por el equipo.}

\subsection{Custodia de evidencias}
Las evidencias originales (carta de aceptaci\'on, consentimientos firmados, audios, fotograf\'ias)
se guardan en una carpeta institucional privada (cuenta \texttt{@uteq.edu.ec}) compartida con el
docente. En el repositorio solo se registra el archivo \texttt{evidencias/MANIFIESTO.md} con el
nombre de archivo, la fecha, la duraci\'on (si aplica) y la huella SHA-256 de cada evidencia, lo
que permite comprobar que el original no fue alterado.

\pendiente{Indicar qui\'en del equipo administra la carpeta, plazo de conservaci\'on y eliminaci\'on.}
